\section{Results}\label{sec:results}

This section reports the main detection results in four steps. Section~\ref{sec:res_single} documents the individual accuracy of each candidate feature and compares them to three trivial firm-pair baselines. Section~\ref{sec:res_composite} reports the logistic composite under both pooled and macro averaging, along with the coefficient estimates. Section~\ref{sec:res_failures} decomposes the fold-level AUCs and characterizes the cartel types that the composite detects reliably. Section~\ref{sec:res_inference_main} reports the cluster-bootstrap confidence intervals and the label-permutation $p$-value.

All estimates use the pair-level analysis sample described in Section~\ref{sec:data_sample}: 100 positive pairs from the six CADE cartels and approximately 100{,}000 negative pairs from non-cartel BEC bidders. The preferred features are the two defined in Section~\ref{sec:data_features}: log pair-level worker flow $\mathrm{wf}_{ij}$ and log product-market overlap $\mathrm{ov}_{ij}$.

\subsection{Single-feature detection accuracy}\label{sec:res_single}

We begin by asking how well each candidate feature identifies same-cartel firm pairs on its own, under leave-one-CADE-cartel-out cross-validation. Table~\ref{tab:single_auc} reports the LOO-CV AUC for each feature and for three trivial baselines.

\begin{table}[htbp]
\centering
\caption{Single-feature LOO-CV AUC}\label{tab:single_auc}
\small
\begin{tabular}{lcc}
\toprule
Feature & Pooled AUC & Macro AUC \\
\midrule
\emph{Proposed fingerprints} \\
\quad Pair-level worker flow $\ln(1 + \mathrm{wf}_{ij})$        & 0.640 & 0.686 \\
\quad Log product-market overlap $\ln(1 + |\Pi_i \cap \Pi_j|)$  & 0.856 & 0.734 \\[4pt]
\emph{Alternative labor feature (rejected)} \\
\quad Static labor cosine similarity (2009) & 0.528 & 0.616 \\[4pt]
\emph{Trivial firm-pair baselines} \\
\quad Four-digit CNAE industry match                            & 0.793 & 0.628 \\
\quad Common-auction indicator                                   & 0.785 & 0.566 \\
\quad Common-buyer indicator                                     & 0.635 & 0.675 \\
\bottomrule
\end{tabular}
\vspace{0.5em}
\begin{minipage}{0.95\textwidth}
\scriptsize{\textbf{Notes:} Pooled AUC is the Mann--Whitney U statistic computed on pooled out-of-fold scores across the six leave-one-CADE-cartel-out folds. Macro AUC is the equal-weighted average of the five informative per-cartel fold AUCs (dropping \emph{sacos de lixo}, which has three positive pairs and a numerically unstable fold AUC). The static labor cosine similarity is reported as the alternative labor feature we rejected after establishing that it fails to generalize cross-cartel under LOO-CV. Source: \texttt{03\_analysis/63\_paper\_master.py}.\par}
\end{minipage}
\end{table}

Three observations bear on the overall picture.

\emph{First}, each candidate fingerprint delivers a modestly informative single-feature signal. Pair-level worker flow alone achieves pooled AUC of 0.64 and macro AUC of 0.69 --- meaningful but not dramatic. Log product-market overlap alone achieves pooled AUC of 0.86 and macro AUC of 0.73. On pooled averaging, overlap is the single strongest feature in the table; on macro averaging, it is mid-pack.

\emph{Second}, the static labor cosine similarity --- the alternative labor feature we tested before selecting pair-level worker flow --- performs at essentially random levels in cross-cartel validation (pooled 0.53, macro 0.62). This is a striking and diagnostic finding. The two labor features are conceptually close but empirically opposite: pair-level worker flow measures \emph{movement} between firms, while static cosine similarity measures \emph{similarity} of firms' labor-market profiles at a point in time. The failure of the static measure and the (modest) success of the flow measure, on the same ground truth, is the core empirical justification for the feature selection we made in Section~\ref{sec:data_features}. We return to this comparison in Section~\ref{sec:rob_features}.

\emph{Third}, the trivial baselines are surprisingly strong. CNAE4-match alone yields pooled AUC of 0.79 and macro AUC of 0.63; the co-bid indicator is only slightly weaker. This sets a meaningful bar that the proposed composite must clear to be interesting as a contribution. We take up that comparison directly in the next subsection.

\subsection{Composite classifier}\label{sec:res_composite}

Table~\ref{tab:composite_coef} reports the in-sample coefficient estimates of the two-feature logistic composite on the standardized features.

\begin{table}[htbp]
\centering
\caption{Logistic composite: coefficient estimates (standardized, class-balanced)}\label{tab:composite_coef}
\small
\begin{tabular}{lcc}
\toprule
Coefficient & Estimate & Interpretation \\
\midrule
Intercept                             & $-1.39$ & class-balanced shift \\
Pair-level worker flow $\mathrm{wf}$  & $+0.83$ & dominant feature on standardized scale\\
Log product-market overlap $\mathrm{ov}$ & $+0.46$ & independent positive contribution \\
\bottomrule
\end{tabular}
\vspace{0.5em}
\begin{minipage}{0.90\textwidth}
\scriptsize{\textbf{Notes:} Logistic regression with class-balanced sample weights, $\ell^2$ regularization $\lambda = 0.001$, and features standardized to zero mean, unit variance in-sample. Both features receive positive weight; on a standardized scale, worker flow's coefficient is approximately 1.8 times the product-market overlap coefficient. Neither feature absorbs the other in the joint specification. Source: \texttt{03\_analysis/63\_paper\_master.py}.\par}
\end{minipage}
\end{table}

Both features contribute positively, and the composite is two-dimensional in the sense that neither feature absorbs the other. On a standardized scale, a one-standard-deviation increase in log pair-level worker flow raises the log-odds of same-cartel membership by $+0.83$; a one-standard-deviation increase in log product-market overlap raises the log-odds by $+0.46$. The ratio of $1.8:1$ tells us that worker flow is the dominant signal on a standardized basis --- which is consistent with the raw-scale comparison in which cartel pairs average 74 workers flowing against a non-cartel baseline of 0.003.

Table~\ref{tab:composite_auc} reports the composite LOO-CV AUCs alongside reference specifications.

\begin{table}[htbp]
\centering
\caption{Composite classifier: leave-one-CADE-cartel-out AUC}\label{tab:composite_auc}
\small
\begin{tabular}{lcc}
\toprule
Specification & Pooled AUC & Macro AUC \\
\midrule
Pair-level worker flow alone                       & 0.640 & 0.686 \\
Log product-market overlap alone (reference)       & 0.856 & 0.734 \\[4pt]
\textbf{Two-feature composite (worker flow $+$ overlap)} & \textbf{0.809} & \textbf{0.845} \\
\quad Cluster bootstrap 95\% CI                    & $[0.480, 0.949]$ & $[0.661, 0.956]$ \\[4pt]
\emph{Comparison specifications} \\
\quad Static labor cosine $+$ overlap (rejected)   & 0.908 & 0.815 \\
\quad CNAE match $+$ overlap (trivial-baseline composite) & 0.855 & 0.804 \\
\quad Three-feature: worker flow $+$ CNAE $+$ overlap   & 0.779 & 0.815 \\
\bottomrule
\end{tabular}
\vspace{0.5em}
\begin{minipage}{0.95\textwidth}
\scriptsize{\textbf{Notes:} The main specification is the two-feature composite (worker flow $+$ overlap). Pooled AUC is dominated by the three large cartels (medicamentos, merenda escolar, trens metros), which contribute 88 of 100 positive pairs; macro AUC equally weights the five informative per-cartel fold AUCs and is the preferred cross-cartel generalization metric. The static-labor-cosine composite is reported as the alternative specification we tested before selecting worker flow; Section~\ref{sec:rob_features} documents the reasoning. The cluster bootstrap (500 resamples over the six CADE cartels with replacement) produces the confidence intervals in the main row. Source: \texttt{03\_analysis/63\_paper\_master.py}.\par}
\end{minipage}
\end{table}

Table~\ref{tab:composite_auc} is the paper's headline. Under the preferred macro-averaging metric, the two-feature composite achieves LOO-CV AUC of $\mathbf{0.845}$, beating every reference specification: overlap alone by $+0.111$, the CNAE-match-plus-overlap composite by $+0.041$, and the static-labor-cosine composite (used in an earlier version of this project) by $+0.030$. Under pooled averaging, the overlap-only reference narrowly beats the main composite ($0.856$ vs. $0.809$), reflecting that overlap fits the three large cartels slightly better while the composite has higher variance across folds.

Why does the composite have higher macro AUC but lower pooled AUC than overlap alone? The answer is concentrated in how each specification handles the distributional tails of the fold-AUC distribution. The composite's per-fold AUCs range from 0.66 (aquecedores solares, cafeteria aeroporto) to 1.00 (merenda escolar) --- a wide spread. Pooled averaging weights each fold by its positive-pair count, so the three large cartels (at AUC 0.84, 1.00, 0.91) pull the pooled average down slightly from what overlap alone would produce on the same large cartels. Macro averaging weights each fold equally, so the high-performing small folds get equal weight with the three large folds, pulling the macro average up. The macro reading favors the composite; the pooled reading is roughly a wash.

For detection deployment, macro is the right metric. Each cartel is a distinct detection target. An authority that has detected the three largest cartels on overlap alone wants to know whether the same screen works on the next cartel --- which may be smaller or different in structure --- and the macro metric is the one that captures this. We therefore present the macro AUC of $0.845$ as the headline and the pooled AUC of $0.809$ as a cross-referential sanity check.

The static-labor-cosine composite (the alternative specification that we tested before selecting worker flow) produces pooled AUC of $0.908$ and macro AUC of $0.815$. Its pooled advantage over the preferred composite is real but concentrated in the three large cartels; its macro disadvantage comes from failing to generalize to the two shortest cartels, where its fold AUC collapses to near chance. On macro aggregation --- again, the preferred cross-cartel metric --- worker flow dominates static similarity by $+0.030$, which is what the feature-selection exercise in Section~\ref{sec:rob_features} reports.

Figure~\ref{fig:composite_roc} plots the ROC curves for the main composite, the overlap-only reference, and the CNAE-match-plus-overlap trivial-baseline composite.

\begin{figure}[htbp]
\centering
\caption{Leave-one-CADE-cartel-out ROC: main composite versus reference specifications}\label{fig:composite_roc}
\includegraphics[width=0.75\textwidth]{../../04_figures/paper_master_roc.pdf}
\begin{minipage}{0.85\textwidth}
\vspace{0.3em}
{\scriptsize \emph{Notes:} Pooled out-of-fold ROC curves for three specifications. Black solid: worker flow $+$ overlap composite (main). Blue dashed: overlap alone. Red dotted: CNAE match $+$ overlap composite. Source: \texttt{03\_analysis/63\_paper\_master.py}.\par}
\end{minipage}
\end{figure}

\subsection{Leave-one-cartel-out fold breakdown}\label{sec:res_failures}

The aggregated macro AUC of 0.845 hides substantial heterogeneity across the six CADE cartels. Because the number of informative cartels is the binding precision constraint of our design, we report the fold-by-fold performance of the main composite directly rather than relying on the aggregate number alone.

\begin{table}[htbp]
\centering
\caption{Main composite leave-one-CADE-cartel-out fold AUCs}\label{tab:fold_auc}
\small
\begin{tabular}{lccl}
\toprule
Held-out cartel & $N_+$ & Fold AUC & Status \\
\midrule
merenda escolar    & 15  & 1.000 & perfect on held-out cartel \\
sacos de lixo      &  3  & 0.997 & near-perfect (small $N_+$) \\
trens metros       & 45  & 0.910 & generalizes \\
medicamentos       & 28  & 0.837 & generalizes \\
aquecedores solares &  6  & 0.663 & weakly generalizes \\
cafeteria aeroporto &  3  & 0.662 & weakly generalizes \\
\midrule
\emph{Pooled (micro)}   & 100 & 0.809 & \\
\emph{Macro (equal weight, excluding sacos)} & & 0.845 & \\
\bottomrule
\end{tabular}
\vspace{0.5em}
\begin{minipage}{0.95\textwidth}
\scriptsize{\textbf{Notes:} Per-fold AUCs under leave-one-CADE-cartel-out cross-validation with the two-feature composite. The \emph{sacos de lixo} fold has only three positive pairs; we retain it in the table and in the pooled average but exclude it from the macro average and from the cluster bootstrap because its AUC is numerically unstable. The \emph{cafeteria aeroporto} fold, which was unusable in earlier drafts because of missing 2009 RAIS occupational data, is now evaluable because the worker-flow feature is defined for any firm pair with RAIS presence in any year from 2009--2017. Source: \texttt{03\_analysis/63\_paper\_master.py}.\par}
\end{minipage}
\end{table}

Four observations.

\emph{First}, the composite generalizes strongly on four of the six cartels. The three cartels with the largest positive-pair counts --- medicamentos (28), merenda escolar (15), trens metros (45) --- yield fold AUCs of 0.84, 1.00, and 0.91 respectively. The single-observation sacos-de-lixo fold (3 positive pairs) scores near 1.00, though its small sample makes the fold AUC imprecise. These four cartels together contribute 91 of the 100 positive pairs and account for essentially all of the composite's pooled AUC.

\emph{Second}, two cartels yield weak generalization. Aquecedores solares (6 positive pairs, one-year conduct period 2009--2010) and cafeteria aeroporto (3 positive pairs, one-year conduct period 2014) both produce fold AUCs of approximately 0.66 --- well above chance, but well below what the composite achieves on the four large cartels. Both are the shortest-lived cartels in the sample. Both have small within-cartel pair counts. We read these two fold AUCs as consistent with the framework prediction that structural fingerprints accumulate over time and that one-year cartels have not had the opportunity to leave detectable traces.

\emph{Third}, the pooled and macro AUCs differ by $0.036$ ($0.809$ vs. $0.845$). The gap comes from how each metric aggregates the heterogeneity. The two weak folds (aquecedores, cafeteria) contribute only 9 of 100 positive pairs to the pooled ranking but receive equal weight with the three large cartels in the macro average. A cartel-size-neutral evaluation favors macro; a detection-deployment metric that reflects the population distribution of cartel sizes favors pooled. We present both and prefer macro for cross-cartel generalization claims.

\emph{Fourth}, the composite's per-fold performance is substantially better than the corresponding numbers under the alternative static-labor-cosine specification we tested before selecting worker flow. Under the static-similarity composite, aquecedores solares's fold AUC is approximately 0.55, trens metros's is 0.83, and the macro average is 0.82. Under the preferred worker-flow composite, those numbers are 0.66, 0.91, and 0.85. The improvement is concentrated in the small and large cartels, not in the medium-sized ones. This is the empirical pattern that led us to select worker flow as the primary labor feature.

\subsection{Statistical inference}\label{sec:res_inference_main}

We test the composite against two complementary null hypotheses. The cluster bootstrap over CADE cartels tests whether the composite's point-estimate AUC (and its increment over overlap alone) is distinguishable from zero in the precision of the five-cluster design. The label-permutation test treats the composite's signal as a whole and asks whether its pooled AUC is distinguishable from what a random label assignment would produce.

\begin{table}[htbp]
\centering
\caption{Cluster bootstrap 95\% CIs on the main composite}\label{tab:boot_ci}
\small
\begin{tabular}{lcc}
\toprule
Quantity & Point estimate & 95\% bootstrap CI \\
\midrule
Main composite pooled AUC                              & 0.809 & $[0.480,\ 0.949]$ \\
Main composite macro AUC                               & 0.845 & $[0.661,\ 0.956]$ \\[4pt]
\emph{Incremental contribution over overlap alone:} \\
\quad AUC difference (pooled)                          & $+0.003$ & $[-0.407,\ +0.266]$ \\
\quad AUC difference (macro)                           & $+0.107$ & $[-0.073,\ +0.265]$ \\
\bottomrule
\end{tabular}
\vspace{0.5em}
\begin{minipage}{0.95\textwidth}
\scriptsize{\textbf{Notes:} Cluster bootstrap with $B = 500$ resamples over the six CADE cartels (with replacement). Each bootstrap resample refits the composite and recomputes the LOO-CV AUC. The AUC difference is the main composite minus the overlap-only baseline. Empirical one-sided $p$-values for $\Delta \leq 0$ are $0.506$ for pooled and $0.104$ for macro. Source: \texttt{03\_analysis/63\_paper\_master.py}.\par}
\end{minipage}
\end{table}

The cluster-bootstrap confidence intervals are wide, reflecting the five-effective-cartel validation sample. The 95\% CI on the macro AUC difference is $[-0.073, +0.265]$, which includes zero. The empirical $p$-value for the null that the macro AUC gain is zero or negative is $0.104$. Under a strict 5\% significance threshold, we cannot reject this null. Under a 10\% threshold, we cannot either --- by four-tenths of a percentage point.

This is the paper's central precision limitation, and we do not hide it. The point-estimate magnitude ($+0.107$ macro AUC) is economically meaningful and is consistent across the bootstrap distribution (the mean of the 500 bootstrap difference replicates is $+0.107$). The statistical imprecision is driven by cartel-level variance --- some bootstrap resamples draw \emph{aquecedores solares} twice, and those samples produce much lower AUC gains, pulling the lower tail of the CI below zero. With a larger ground-truth sample (more informative cartels), the CI would tighten; with the five cartels we have, this is the precision that a cluster bootstrap can produce.

The label-permutation test addresses a complementary question. For $B = 500$ permutations of the label vector, we recompute the LOO-CV pooled AUC of the composite on the permuted sample. The null distribution is centered tightly at $0.595$ with a 99th percentile of $0.813$. The observed pooled AUC of $0.809$ is at the 98th percentile of the null, giving an empirical $p$-value of $\mathbf{0.020}$. We reject at the 5\% level the hypothesis that the composite's pooled signal is a label-assignment artifact.

This pair of inference tests answers different questions about the same data, and we report both deliberately. The \emph{absolute-signal test} --- cluster-aware label permutation --- asks whether the composite's pooled ranking is distinguishable from what a random label assignment would produce. Against that null, we reject at the 2\% level: the composite has real signal. The \emph{incremental-contribution test} --- cluster bootstrap over cartels on the AUC \emph{difference} between the two-feature composite and the overlap-only baseline --- asks whether a specific feature's marginal effect is distinguishable from zero. Against that null, at the 5\% threshold, we cannot reject. The two tests are not substitutes for each other, and neither is dispositive on its own.

We emphasize this distinction because applied detection work has historically reported one or the other without naming the difference. A paper that shows its classifier beats random labels is not the same as a paper that shows its classifier beats a simpler reference classifier; the two claims live in different statistical worlds, and collapsing them invites either overclaiming (reporting only the easier test) or underclaiming (reporting only the harder one). In a setting like ours --- where the ground truth is small, the cluster bootstrap inherits the small-cluster variance, and the point-estimate magnitude is economically meaningful but imprecise --- the honest summary requires both tests and refuses to reduce them to a single number.

A reader who weighs point-estimate magnitude, theoretical grounding, and the absolute-signal test alongside the incremental-contribution bootstrap will conclude that the composite captures a real detection signal and that worker flow contributes to it in an economically meaningful way. A reader who applies a strict 5\% incremental-contribution threshold will not reach that conclusion on the statistics alone. Both readings are legitimate. We invite future detection work to report both tests and to name the distinction, so that claims of the form ``our classifier works'' and ``our feature adds something'' can be evaluated on their own terms.

\medskip

\noindent\textbf{Summary.} Under leave-one-CADE-cartel-out cross-validation, the two-feature composite of pair-level worker flow and log product-market overlap achieves a macro AUC of $0.845$ and a pooled AUC of $0.809$ on the six-cartel validation set. The macro AUC beats every comparison specification: overlap alone by $+0.111$, CNAE-match-plus-overlap by $+0.041$, and the static-labor-cosine composite by $+0.030$. The composite generalizes on four cartels at fold AUCs between 0.84 and 1.00 and weakly generalizes on the two shortest-lived cartels at fold AUCs of approximately 0.66. Cluster bootstrap over cartels gives a macro AUC 95\% CI of $[0.661, 0.956]$ and a 95\% CI on the incremental contribution of worker flow over overlap alone of $[-0.073, +0.265]$ with empirical $p$-value $0.104$. Label permutation on the pooled AUC rejects the no-signal null at the 2\% level. Section~\ref{sec:robustness} reports further sensitivity checks.
